\documentclass[12pt]{article}
\usepackage[T1]{fontenc}
\usepackage[utf8]{inputenc}
\usepackage[brazil]{babel}
\usepackage[margin=1in]{geometry}
\setlength{\parindent}{0pt}
\begin{document}
\section*{Tema 12}
Processo: Processo\\
Situacao: Revisado - Tema 216/TNU\\
Ramo: DIREITO PREVIDENCIÁRIO\\
Questao: Saber se pagamento de remuneração indireta a aluno aprendiz autoriza respectiva contagem de tempo de serviço para fins previdenciários.\\
Tese: Para fins previdenciários, o cômputo do tempo de serviço prestado como aluno-aprendiz exige a comprovação de que, durante o período de aprendizado, houve simultaneamente: (i) retribuição consubstanciada em prestação pecuniária ou em auxílios materiais; (ii) à conta do Orçamento; (iii) a título de contraprestação por labor; (iv) na execução de bens e serviços destinados a terceiros. (Tese firmada no Tema 216/TNU). Vide Súmula 18 da TNU.\\
Fonte: NAO\_INFORMADO
\end{document}
